\begin{table}[!htbp]\centering
\caption{Closure-catchment influence diagnostics}
\label{tab:closure-level-inference}
\small
\resizebox{\textwidth}{!}{%
\begin{tabular}{lccccccc}
\toprule
Outcome & Baseline ATT & Leave-one-out range & Max $|\Delta|$ & Most influential closure & Year & Omitted munis & Any CI excludes 0 \\
\midrule
Suicide & 0.28 & [0.02, 0.59] & 0.31 & 5585279\_2012 & 2012 & 5 & no \\
Self-harm & 0.61 & [-0.41, 1.00] & 1.02 & 2291061\_2012 & 2012 & 5 & no \\
\bottomrule
\multicolumn{8}{p{0.96\textwidth}}{\footnotesize Notes: The diagnostic re-estimates the population-weighted Sun--Abraham headline specification after omitting, one PNASH closure at a time, all municipalities exposed to that closure by the patient-flow rule ($\theta=0.05$, all SUS admissions in year $-1$). Exact closure-catchment clustering is not well-defined in this panel because some municipalities can be exposed to more than one closing hospital and never-treated controls have no closure catchment. The table therefore reports leave-one-closure-catchment-out influence rather than a clustered standard error at the closure level.}\\
\end{tabular}
}%
\end{table}
